\section{Density of the original joint boundaries and determinant dissipation}
\label{sec:joint-density-dissipation}

Retain the complete packet \(h\), its degree \(d\ge1\), the
arithmetic unit \(\upsilon_h\), and every original map in
Section~\ref{sec:coherent-tensor-integration}.  In particular
\[
 s_j=\tfrac12+it_j,\quad v_h=g/h,\quad g=2\xi,\quad
 d\nu_h(t)=|v_h(\tfrac12+it)|^2\frac{dt}{2\pi},\quad
 E^{(k)}=E_h^{\otimes k},\quad \mathcal H_k=L^2(d^kx).
 \tag{JD.1}
\]
The one-variable cyclic-density argument in the coherent-interpolation
source has a full joint-source extension.  We prove it here, including
the exact relation between the disappearing Hilbert representatives
and their unchanged finite jets.

\subsection{Polynomial density with the original weighted measures}

\begin{lemma}
Let \(\mu\) be a finite positive Borel measure on \(\mathbb R^k\)
such that \(\int e^{a\sum_j|t_j|}\,d\mu<\infty\) for some
\(a>0\).  The complex polynomials in \(t_1,\ldots,t_k\) are
dense in \(L^2(\mu)\).
\end{lemma}
\begin{proof}
Suppose \(f\in L^2(\mu)\) is orthogonal to every polynomial.
Define the finite complex measure \(d\eta=\overline f\,d\mu\).
For \(0<b<a/2\), Cauchy--Schwarz gives
\[
 \int e^{b\sum_j|t_j|}\,d|\eta|
 \le \|f\|_{L^2(\mu)}
       \left(\int e^{2b\sum_j|t_j|}\,d\mu\right)^{1/2}<\infty.
 \tag{JD.2}
\]
Thus \(\Phi(z)=\int e^{i\sum_jz_jt_j}\,d\eta(t)\) is
holomorphic in the connected tube \(|\operatorname{Im}z_j|<b\).
On every smaller closed tube, each differentiated integrand is
dominated by the integrable function in (JD.2), multiplied by a
constant; the strict exponent margin absorbs each monomial.  This
justifies every derivative and the convergent Taylor series on small
polydiscs.  All derivatives at zero vanish by polynomial
orthogonality.  The Taylor series therefore vanishes near zero, and
successive overlapping polydiscs along any path in the tube show
\(\Phi=0\) throughout it.  In particular the Fourier transform
of \(\eta\) vanishes on \(\mathbb R^k\).

For clarity, uniqueness here follows from an explicit approximate
identity.  Put
\[
 \gamma_\epsilon(x)=(4\pi\epsilon)^{-k/2}
                   e^{-|x|^2/(4\epsilon)}
   =(2\pi)^{-k}\int e^{-\epsilon|u|^2}e^{iu\cdot x}\,du.
 \tag{JD.3}
\]
The second equality is the product Gaussian integral, obtained in
one variable by differentiating its Fourier integral in \(x\),
integrating by parts in \(u\), and evaluating its value at zero.
The Gaussian integrability and finiteness of \(|\eta|\) justify
Fubini.  Consequently \(\eta*\gamma_\epsilon=0\).  For any
continuous compactly supported \(\varphi\), integration and
Fubini give \(\int(\varphi*\gamma_\epsilon)\,d\eta=0\).
These convolutions converge uniformly to \(\varphi\): split the
integral at a fixed small radius, use uniform continuity within it,
and use the Gaussian tail and boundedness outside it.  Hence
\(\int\varphi\,d\eta=0\).  Uniqueness of finite Borel measures
from their continuous compactly supported integrals gives
\(\eta=0\), and therefore \(f=0\) almost everywhere.  The
orthogonal complement of the polynomial closure is zero, proving
the lemma.
\end{proof}

The actual measure \(\nu_h\) has every exponential moment of
order less than \(\pi/2\), with its full factor \(dt/(2\pi)\),
as proved in (GC.2).  Multiplying it by any fixed polynomial in
\(t\) preserves the same open range of exponential moments:
choose a strictly larger exponent still below \(\pi/2\) and
absorb the polynomial into the difference.  Product integration
therefore verifies the lemma for both measures
\[
 \nu_h^{\otimes k},\qquad
 d\mu_{h,k}(t)=|h(s_1)|^2\prod_{j=1}^k d\nu_h(t_j).
 \tag{JD.4}
\]
No factor \(h\) has been discarded.  Its finitely many zeros on
the first integration line have measure zero.  The multiplication
map
\[
 U_h:L^2(\mu_{h,k})\longrightarrow L^2(\nu_h^{\otimes k}),
 \qquad U_hP=h(s_1)P
 \tag{JD.5}
\]
is a unitary surjection: its norm identity follows directly from
(JD.4), and its inverse is division by \(h(s_1)\) away from
that null set, with an arbitrary value on the null set.  Polynomials
in \(s_j=1/2+it_j\) are precisely polynomials in \(t_j\), since
these coordinate changes are invertible.  Thus (JD.5) and the lemma
prove density of \(h(s_1)\mathbb C[s_1,\ldots,s_k]\) in
\(L^2(\nu_h^{\otimes k})\).  This dense space is contained in
the original joint ideal \((h(s_1),\ldots,h(s_k))\).

\subsection{Strong limits of the actual source representatives}

The Mellin Plancherel map, with its original sign and measure, is
\[
 (\mathcal M_kF)(t)=\int_{(0,\infty)^k}
    F(x)\prod_{j=1}^k x_j^{-1/2+it_j}\,d^kx,
 \quad
 \mathcal M_k:\mathcal H_k\longrightarrow
                 L^2\left(\frac{d^kt}{(2\pi)^k}\right).
 \tag{JD.6}
\]
Indeed the substitutions \(x_j=e^{u_j}\) send \(F\) to
\(e^{\sum_j u_j/2}F(e^{u_1},\ldots,e^{u_k})\) isometrically
in \(L^2(d^ku)\), and (JD.6) becomes its Fourier transform
with kernel \(e^{it\cdot u}\).  The inverse has kernel
\(e^{-it\cdot u}\) and factor \((2\pi)^{-k}\).
Multiplication by \(\prod_jv_h(s_j)\) is a unitary surjection
from \(L^2(\nu_h^{\otimes k})\) to the target in (JD.6).
For surjectivity divide by this product almost everywhere: the
entire nonzero function \(v_h\) has only isolated zeros on
the line, and the product zero set is null.  Composing with the
inverse of (JD.6) gives the unitary extension
\[
 \mathcal U_{h,k}:L^2(\nu_h^{\otimes k})\longrightarrow
 \mathcal H_k,\qquad
 \mathcal U_{h,k}P=\mathcal T_h^{(k)}P
                 \quad\hbox{for every polynomial }P.
 \tag{JD.7}
\]
The polynomial identity follows from the actual Mellin transform of
\(F_h\) and \(D_j\), so this is an extension of the original
source map, not a change of representative.

\begin{theorem}[Joint boundary density and exact finite-jet graph closure]
For each \(k\ge1\), the original total-degree boundary spaces
\(\mathcal B_M\) satisfy
\[
 \overline{\bigcup_{M\ge k(d-1)}\mathcal B_M}=\mathcal H_k,
 \qquad P_{\mathcal B_M}\longrightarrow I\quad\hbox{strongly}.
 \tag{JD.8}
\]
For the actual coherent minima of (CT.13), in the original fixed
finite monomial tensor coordinates,
\[
 \|R_M\|_{E^{(k)}\to\mathcal H_k}\longrightarrow0,
 \qquad G_M\downarrow0,
 \qquad \det G_M\longrightarrow0.
 \tag{JD.9}
\]
Here \(G_M>0\) at every finite \(M\), and the middle limit
is both decreasing in the positive-semidefinite order and convergent
in matrix operator norm.  Every finite jet is still retained:
\(J^{(k)}R_M=I\) for every \(M\).

Define \(J_{\rm pol}\) on the dense original source
\(\mathcal D_{h,k}=\mathcal T_h^{(k)}\mathbb C[s_1,\ldots,s_k]\)
by its complete arithmetic jets, including multiplication by
\(\upsilon_h^{\otimes k}\).  Its graph has exactly the closure
\[
 \overline{\operatorname{Graph}(J_{\rm pol})}
       =\mathcal H_k\oplus E^{(k)}.
 \tag{JD.10}
\]
As a closed linear relation this has the split exact sequence
\[
 0\longrightarrow E^{(k)}
   \xrightarrow{\ u\mapsto(0,u)\ }
 \overline{\operatorname{Graph}(J_{\rm pol})}
   \xrightarrow{\ (F,u)\mapsto F\ }\mathcal H_k
   \longrightarrow0,
 \quad F\longmapsto(F,0).
 \tag{JD.11}
\]
\end{theorem}
\begin{proof}
Every polynomial in \(h(s_1)\mathbb C[s_1,\ldots,s_k]\)
belongs to one finite total-degree ideal \(\mathcal I_{h,k,M}\).
Its density proved in (JD.5), followed by the unitary map (JD.7),
proves the first assertion in (JD.8).  Given \(F\in\mathcal H_k\)
and \(\epsilon>0\), choose \(b\in\mathcal B_N\) with
\(\|F-b\|<\epsilon\).  For every \(M\ge N\), the distance
from \(F\) to \(\mathcal B_M\) is at most that norm, so
\(\|(I-P_{\mathcal B_M})F\|<\epsilon\).  This proves strong
convergence.

Apply it to each of the \(d^k\) images of the fixed monomial
basis under \(s_h^{\otimes k}\).  The identity
\(R_M=(I-P_{\mathcal B_M})s_h^{\otimes k}\) makes each
column tend to zero.  For a vector with coordinate norm one,
Cauchy--Schwarz bounds its image by the square root of the sum of
the squared column norms, proving the first limit in (JD.9).
Since \(G_M=R_M^*R_M\) and
\(G_M-G_{M+1}=Y_M^*Y_M\), the norm limit and monotonicity
follow.  The determinant limit follows by continuity of the
determinant in these same finite coordinates.  Positivity and
\(J^{(k)}R_M=I\) are the proved finite interpolation identities;
none of these finite algebraic equations is lost in taking the
Hilbert norm limit.

The kernel of \(J_{\rm pol}\) is exactly the union of the
\(\mathcal B_M\), so its graph closure contains
\(\mathcal H_k\oplus\{0\}\) by (JD.8).  For any
\(u\in E^{(k)}\), the vectors \((R_Mu,u)\) belong to its
graph and converge to \((0,u)\) by (JD.9).  The closure is a
linear subspace, and therefore contains every \((F,u)\); the
reverse inclusion is its declared ambient space.  This proves
(JD.10), and the maps, exactness and section in (JD.11) follow
immediately by evaluating their two coordinates.  In particular
the closure is a relation with vertical fibre \(E^{(k)}\),
and it cannot be the graph of a single-valued operator because
\(d^k>0\).
\end{proof}

There are corresponding exact limits for the original weight form
and both interpolation kernels.  Write \(U=U_{\upsilon_h^{\otimes k}}\),
let \(K_M\) be the remainder-only polynomial kernel on the joint
source, and let \(\mathcal C_M=G_M^{-1}=UK_MU^*\) be the
complete arithmetic interpolation matrix.  In the same fixed
coordinates,
\[
 \begin{split}
 \|W_M\|&\le(2\|A^{(k)}\|+k)\|G_M\|\longrightarrow0,\\
 \lambda_{\min}(\mathcal C_M)&=\|G_M\|^{-1}
                                     \longrightarrow\infty,\\
 \lambda_{\min}(K_M)&\ge
                \frac{1}{\|U\|^2\|G_M\|}\longrightarrow\infty.
 \end{split}
 \tag{JD.9a}
\]
The first line follows by the triangle inequality applied to the
definition of \(W_M\).  Positivity of \(G_M\) proves the
second.  For the third, use
\(K_M=U^{-1}\mathcal C_MU^{-*}\) and
\(\|U^{-*}z\|\ge\|z\|/\|U\|\).  Thus these coercive
limits include every retained local jet, and the arithmetic unit
remains in the comparison.  The relative weight is related to its
absolute form by the exact congruence
\[
 W_M=G_M^{1/2}S_MG_M^{1/2},\qquad
 \lambda_{\min}(G_M)\epsilon_M
    \le\|W_M\|\le\|G_M\|\epsilon_M.
 \tag{JD.9b}
\]
The upper bound is submultiplicativity; applying it instead to
\(S_M=G_M^{-1/2}W_MG_M^{-1/2}\) gives the lower bound.
Equations (JD.9a)--(JD.9b) retain the metric factors needed to
pass from the proved absolute limits to the relative quantity.

The complementary quotient of the same closed relation is also
explicit:
\[
 \frac{\overline{\operatorname{Graph}(J_{\rm pol})}}
            {\mathcal H_k\oplus\{0\}}
 \xrightarrow{\ \sim\ }E^{(k)},\qquad
 [(F,u)]\longmapsto u,\qquad
 u\longmapsto[(0,u)].
 \tag{JD.11a}
\]
The kernel, well-definedness and both inverse identities follow by
subtracting the two displayed pairs.  With the product Hilbert norm
the quotient norm is exactly \(\|u\|\), since
\(\inf_{H\in\mathcal H_k}(\|F-H\|^2+\|u\|^2)^{1/2}=\|u\|\),
attained at \(H=F\).  Thus every completed-boundary class retains
its original finite jet through an explicit isometric quotient.

This limit has an exact Split-Zero morphism.  For a complex vector
space \(X\), use its one-support lift
\(\widehat X=\{\tau_X\}\sqcup X^\bullet\), with
\(\tau_X\) the additive identity, vector addition within
\(X^\bullet\), and the original \(G(\mathbb C)\)-scalar
action.  In particular \(0_X^\bullet\ne\tau_X\).
Any linear map \(T:X\to Y\) lifts to
\(\widehat T(\tau_X)=\tau_Y\) and
\(\widehat T(u^\bullet)=(Tu)^\bullet\); direct substitution
checks addition and every scalar, including the supported scalar
zero and external scalar zero.  Give the supported fibre its norm
topology and \(\tau_X\) its isolated topology.  Equation (JD.9)
then says exactly
\[
 \widehat R_M\longrightarrow\widehat 0,
 \qquad \widehat0(\tau_E)=\tau_{\mathcal H},\quad
 \widehat0(u^\bullet)=0_{\mathcal H}^\bullet
                    \quad(u\in E^{(k)}).
 \tag{JD.12}
\]
Thus the limiting morphism retains the support of every input.
Equation (JD.11) identifies all finite jets in the vertical fibre
of the closed relation; it supplies the full typed relation between
these supported Hilbert limits and the unchanged algebraic jets.

\subsection{The accumulated determinant cost of the same boundary sequence}

Continue to use the original ordered spectral sum and layer energy
\[
 \mathcal D_k=\sum_{i_1,\ldots,i_k}
     \left|\sum_{j=1}^k(\operatorname{Re}\rho_{i_j}-1/2)\right|,
 \qquad \chi_M=\operatorname{Tr}(G_M^{-1}C_M^*C_M),
 \quad \tau_M=\operatorname{Tr}(G_M^{-1}(G_M-G_{M+1})).
 \tag{JD.13}
\]
The list \(\rho_1,\ldots,\rho_d\) includes every algebraic
multiplicity.  Define \(\pi_M=\det G_{M+1}/\det G_M\).
For integers \(k(d-1)\le a<b\), the following exact finite
inequality holds whenever \(\mathcal D_k>0\):
\[
 \boxed{
 \mathcal D_k^2\sum_{M=a}^{b-1}\frac1{\chi_M}
 \le \sum_{M=a}^{b-1}\tau_M
 \le \log\frac{\det G_a}{\det G_b}.}
 \tag{JD.14}
\]
All denominators in this statement are proved positive.  Indeed
(CT.24) gives \(\mathcal D_k^2\le\tau_M\chi_M\), with
both factors nonnegative; strict positivity of the left side
forces strict positivity of each factor.  Let
\(t_{M,j}\in[0,1)\) be the full eigenvalue list of
\(G_M^{-1/2}(G_M-G_{M+1})G_M^{-1/2}\), including zeros.
The finite original matrices give exactly
\[
 \tau_M=\sum_jt_{M,j},\quad
 \pi_M=\prod_j(1-t_{M,j}),\quad
 -\log\pi_M-\tau_M
   =\sum_j\int_0^{t_{M,j}}\frac{u}{1-u}\,du\ge0.
 \tag{JD.15}
\]
The first two identities are congruence and determinant identities;
the last follows by integrating the derivative of
\(-\log(1-t)-t\).  Sum (CT.24) divided by \(\chi_M\),
then sum (JD.15).  The product of the determinant ratios telescopes
without a change of coordinates, proving (JD.14).

If \(\mathcal D_k=0\), every repeated tuple
\((i,\ldots,i)\) has zero summand, hence
\(\operatorname{Re}\rho_i=1/2\) for every retained zero.
Conversely these equalities give \(\mathcal D_k=0\).
The second inequality in (JD.14) and the exact remainder (JD.15)
hold also in this case; no quotient \(0/0\) is introduced.
Equation (JD.9) proves that the right-hand side tends to infinity
as \(b\to\infty\) for each fixed \(a\).  Equations
(JD.13)--(JD.15) retain the actual rate through \(\chi_M\),
the complete relative Gram losses and their explicit logarithmic
remainder.  They impose no unproved upper estimate on that rate.

The same sequence supplies two further exact energy identities:
\[
 G_a=\sum_{M=a}^\infty Y_M^*Y_M,\qquad
 \sum_{M=a}^\infty\|Y_M\|_{\mathrm{HS}}^2=\operatorname{Tr}G_a,
 \qquad
 \sum_{M=a}^\infty\tau_M=\infty.
 \tag{JD.16}
\]
The first series converges in matrix operator norm, by telescoping
\(G_M-G_{M+1}=Y_M^*Y_M\) and the proved limit \(G_M\to0\).
Taking traces of its finite partial sums proves the second identity;
all summands are nonnegative.  To prove the third, suppose instead
that \(\sum_M\tau_M\) were finite.  Then beyond some index all
\(\tau_M\le1/2\), so every \(0\le t_{M,j}\le1/2\).
Integrating \((1-u)^{-1}\le2\) over \([0,t_{M,j}]\) gives
\(-\log(1-t_{M,j})\le2t_{M,j}\).  Hence the tail sum of
\(-\log\pi_M\) would be bounded by \(2\sum_M\tau_M<\infty\).
The finite initial terms are finite because \(G_M>0\).
Telescoping would therefore bound
\(\log(\det G_a/\det G_b)\) uniformly in \(b\), contradicting
\(\det G_b\to0\).  This proves the divergence, including
packets entirely on the critical line and every retained
nilpotent chain.  The equality
\(\tau_M=\|Y_MG_M^{-1/2}\|_{\mathrm{HS}}^2\) identifies
the relative energy in this divergence with the same layer map
appearing in the convergent absolute energy series.
